\chapter{Wittig e outras reações que formam C=C}\label{ch:b2:wittig}

A fêmea da mosca doméstica atrai os machos com um hidrocarboneto de 23 carbonos que leva uma única
ligação dupla, entre os carbonos 9 e 10, de geometria (\textit{Z}). Para prepará-lo em quantidade, o
químico precisa colocar essa ligação dupla exatamente ali, e com essa geometria: uma eliminação a
espalharia por várias posições e daria as duas geometrias. Este capítulo constrói as ligações duplas
pelo outro caminho, unindo dois pedaços em um grupo carbonila, de modo que a ligação dupla aparece
onde estava a \ce{C=O}.
% ledger: use:muscalure

\begin{recall}
\Cref{ch:b2:enolates-aldol}: a \omterm{def:b2:enolates-aldol:aldol}{condensação aldólica}. \Cref{ch:b2:alkene-redox}: a hidrogenação de um
alcino a alceno (\textit{Z}) sobre um catalisador envenenado. O volume do 1.º ano: E1, E2 e a regra de
Zaitsev, a substituição bimolecular, os reagentes de Grignard e a inversão de polaridade, os
descritores (\textit{Z})/(\textit{E}). O volume escolar: a economia atômica.
\end{recall}

\section{Sais de fosfônio e ilídeos}

\begin{definition}[Ilídeos de fosfônio]\label{def:b2:wittig:ylide}
Um \emph{sal de fosfônio}\index{sal de fosfonio@sal de fosfônio} $\mathrm{R_3P^+{-}CH_2R'\ X^-}$ é preparado por substituição
bimolecular de um haloalcano por uma fosfina, em geral a trifenilfosfina \ce{Ph3P}. Retirar um
hidrogênio do carbono ligado ao fósforo dá um \emph{ilídeo de fosfônio}\index{ilideo de fosfonio@ilídeo de fosfônio},
uma molécula neutra com cargas opostas adjacentes, $\mathrm{Ph_3P^+{-}CH^-{-}R'}$, também escrita $\mathrm{Ph_3P{=}CHR'}$. Um
\emph{ilídeo estabilizado}\index{ilideo estabilizado@ilídeo estabilizado} leva nesse carbono um grupo retirador de elétrons
(um éster, uma cetona, uma \omterm{def:b2:amines:nitrile}{nitrila}) que deslocaliza a carga negativa.
\end{definition}

\begin{omfigure}
\setchemfig{atom sep=1.9em}
\schemestart
\chemfig{Ph_3P} \+ \chemfig{CH_3-I}
\arrow{->[S$_\mathrm{N}$2]}[,0.9]
\chemfig{Ph_3P^{+}-CH_3}\;\chemfig{I^{-}}
\arrow{->[\ce{BuLi}][$-$ \ce{BuH}, \ce{LiI}]}[,1.3]
\chemfig{Ph_3P^{+}-\chemabove{C}{\scriptstyle -}H_2}
\arrow{<->}[,0.8]
\chemfig{Ph_3P=CH_2}
\schemestop
\par\vspace{6pt}

{\small O metilidenotrifenilfosforano, o ilídeo mais simples. A trifenilfosfina desloca o iodeto do
iodometano; o butil-lítio retira um próton do grupo metila. As duas estruturas à direita são duas
maneiras de escrever a mesma molécula: a de cargas separadas é a mais próxima da realidade.}
\end{omfigure}

\begin{proposition}[Preparar o ilídeo]\label{prop:b2:wittig:ylide-acidity}
Um \omterm{def:b2:wittig:ylide}{sal de fosfônio} é muito mais ácido no carbono vizinho do fósforo que um alcano. Um sal de
alquilfosfônio simples ainda exige uma base muito forte (butil-lítio, hidreto de sódio, amideto de sódio)
em condições anidras e inertes; um sal cujo \ce{CH2} leva também um grupo carbonila é desprotonado por
uma base fraca (carbonato de sódio, hidróxido de sódio), e o \omterm{def:b2:wittig:ylide}{ilídeo estabilizado} formado pode ser
isolado e armazenado.
\end{proposition}

\begin{proof}[Raciocínio]
O fósforo positivo vizinho do carbânion o estabiliza por sua carga e pela polarizabilidade do grande
átomo de fósforo: um ilídeo simples é um carbânion estabilizado, mas menos que um \omterm{def:b2:enolates-aldol:enolate}{enolato}, daí a base
forte. Uma carbonila adjacente espalha a carga negativa sobre o oxigênio, como em um \omterm{def:b2:enolates-aldol:enolate}{enolato}
(\cref{prop:b2:enolates-aldol:alpha-acidity}), além do fósforo: o ácido conjugado torna-se ácido o
bastante para uma base fraca, e o ilídeo, muito menos básico, já não reage com a água nem com o ar.
\end{proof}

\section{A reação de Wittig}

\begin{definition}[Reação de Wittig]\label{def:b2:wittig:wittig}
A \emph{reação de Wittig}\index{reacao de Wittig@reação de Wittig} de um \omterm{def:b2:wittig:ylide}{ilídeo de fosfônio} com um aldeído ou uma cetona
dá um alceno e o óxido de trifenilfosfina:
\begin{equation*}
\mathrm{Ph_3P{=}CHR + R'CHO \longrightarrow R'CH{=}CHR + Ph_3P{=}O}.
\end{equation*}
Ela passa por um \emph{oxafosfetano}\index{oxafosfetano}, um anel de quatro membros formado pelo fósforo, dois
carbonos e o oxigênio, criado quando o carbono do ilídeo se liga ao carbono carbonílico enquanto o
oxigênio carbonílico se liga ao fósforo.
\end{definition}

\begin{omfigure}
\setchemfig{atom sep=1.9em}
\schemestart
\chemfig{Ph_3P=CH-R} \+ \chemfig{O=CH-R'}
\arrow{->[adição]}[,1.0]
\chemfig{Ph_3P?-[6]CH(-[4]R)-[0]CH(-[0]R')-[2]O?}
\schemestop

\medskip
\schemestart
\arrow{->[eliminação]}[,1.1]
\chemfig{R-CH=CH-R'} \+ \chemfig{Ph_3P=O}
\schemestop
\par\vspace{6pt}

{\small A \omterm{def:b2:wittig:wittig}{reação de Wittig}. O carbono do ilídeo e o carbono carbonílico se ligam enquanto o oxigênio se
liga ao fósforo, em uma única etapa: o \omterm{def:b2:wittig:wittig}{oxafosfetano}. Seu anel se abre então no outro sentido, rompendo
as ligações \ce{P-C} e \ce{C-O}: a ligação \ce{C=C} se forma entre os dois antigos carbonos reativos, e
a ligação \ce{P=O} do óxido de trifenilfosfina se forma ao mesmo tempo.}
\end{omfigure}

\begin{proposition}[Uma ligação dupla exatamente no lugar]\label{prop:b2:wittig:regiospecific}
Em uma \omterm{def:b2:wittig:wittig}{reação de Wittig}, a nova ligação \ce{C=C} une o antigo carbono carbonílico ao antigo carbono do
ilídeo, e em nenhum outro lugar.
\end{proposition}

\begin{proof}[Raciocínio]
As únicas ligações formadas e rompidas são as do anel de quatro membros: o carbono do ilídeo com o
carbono carbonílico, depois as ligações \ce{C-P} e \ce{C-O}. Nenhum hidrogênio se desloca e nenhum
carbocátion ou carbânion se forma ao longo da cadeia, de modo que a ligação dupla não pode migrar. Uma
eliminação, ao contrário, escolhe entre os hidrogênios de vários carbonos vizinhos (regra de Zaitsev),
e uma reação E1 pode sofrer rearranjo.
\end{proof}

\begin{proposition}[Força motriz]\label{prop:b2:wittig:driving}
A formação da ligação \ce{P=O} do óxido de trifenilfosfina torna a \omterm{def:b2:wittig:wittig}{reação de Wittig} fortemente
favorável e irreversível.
\end{proposition}

\begin{proof}[Raciocínio]
A reação troca uma ligação \ce{C=O} e uma ligação \ce{P-C} (no ilídeo, a \ce{P=C} é sobretudo uma
ligação simples \ce{P+-C-}) por uma ligação \ce{C=C} e uma ligação \ce{P=O}. A ligação \ce{P=O} de um
óxido de fosfina é muito forte, muito mais forte que a ligação \ce{P-C} perdida, e a \ce{C=O} perdida é
compensada em parte pela \ce{C=C} formada: o balanço é claramente descendente. O fósforo é o
sorvedouro de oxigênio da reação.
\end{proof}

\section{Estereosseletividade}

\begin{proposition}[Geometria do alceno]\label{prop:b2:wittig:stereo}
Com um aldeído, um ilídeo não estabilizado (\ce{Ph3P=CHR}, R um grupo alquila) dá sobretudo o alceno (\textit{Z}),
em condições isentas de sais. Um \omterm{def:b2:wittig:ylide}{ilídeo estabilizado} (\ce{Ph3P=CH-COOR}, \ce{Ph3P=CH-COR}) dá sobretudo o
alceno (\textit{E}). Um \omterm{def:b2:wittig:ylide}{ilídeo estabilizado} apenas por um grupo arila dá misturas.
\end{proposition}

\admitted

São observações. Sua explicação habitual lê a \omterm{def:b2:wittig:wittig}{reação de Wittig} como uma competição entre velocidades e
estabilidades. Com um ilídeo reativo, não estabilizado, o \omterm{def:b2:wittig:wittig}{oxafosfetano} se forma de modo rápido e
irreversível, pelo estado de transição menos congestionado, no qual os dois substituintes acabam do
mesmo lado do anel; ele depois se desfaz conservando esse arranjo: o controle cinético dá o
alceno (\textit{Z}). Com um \omterm{def:b2:wittig:ylide}{ilídeo estabilizado}, a primeira etapa é mais lenta e reversível; os
\omterm{def:b2:wittig:wittig}{oxafosfetanos} se interconvertem passando pelos reagentes de partida, e aquele que tem os substituintes
em lados opostos, menos tensionado, vence: o alceno (\textit{E}), sob controle termodinâmico.

\begin{omfigure}
\setchemfig{atom sep=1.9em}
\schemestart
\chemfig{Ph_3P=CH-CH_2CH_3} \+ \chemfig{O=CH-Ph}
\arrow{->}[,0.8]
\chemfig{Ph-[:-30]=[:30]-[:90]CH_2CH_3}
\schemestop

\medskip
\schemestart
\chemfig{Ph_3P=CH-COOEt} \+ \chemfig{O=CH-Ph}
\arrow{->}[,0.8]
\chemfig{Ph-[:-30]=[:30]-[:-30]COOEt}
\schemestop
\par\vspace{6pt}

{\small Dois ilídeos, um aldeído. Em cima: o ilídeo propilideno, não estabilizado, dá sobretudo o
(\textit{Z})-1-fenilbut-1-eno. Embaixo: o \omterm{def:b2:wittig:ylide}{ilídeo estabilizado} dá sobretudo o
(\textit{E})-3-fenilpropenoato de etila (cinamato de etila). O óxido de trifenilfosfina, formado nos
dois casos, não está representado.}
\end{omfigure}

\section{A reação de Horner--Wadsworth--Emmons}

\begin{definition}[Fosfonatos e a reação HWE]\label{def:b2:wittig:hwe}
Um \emph{fosfonato}\index{fosfonato} é um éster de um ácido fosfônico, $\mathrm{(RO)_2P({=}O){-}R'}$, com
uma ligação \ce{P-C}. Ele é preparado pela reação de Michaelis--Arbuzov de um fosfito de trialquila com
um haloalcano, $\mathrm{P(OEt)_3 + R'Br \to (EtO)_2P(O)R' + EtBr}$. Na
\emph{reação de Horner--Wadsworth--Emmons}\index{reacao de Horner--Wadsworth--Emmons@reação de Horner--Wadsworth--Emmons} (HWE), o ânion de um
fosfonato que leva um grupo retirador de elétrons em seu \ce{CH2} se adiciona a um aldeído e dá um
alceno e um ânion fosfato de dialquila.
\end{definition}

\begin{omfigure}
\setchemfig{atom sep=1.8em}
\schemestart
\chemfig{EtO-P(=[2]O)(-[6]OEt)-CH_2-COOEt}
\arrow{->[\ce{NaH}][$-$ \ce{H2}]}[,0.9]
\chemfig{EtO-P(=[2]O)(-[6]OEt)-\chemabove{C}{\scriptstyle -}H-COOEt}
\schemestop

\medskip
\schemestart
\arrow{->[\ce{PhCHO}]}[,0.9]
\chemfig{*6(=-=(-[:30]=[:-30]-[:30](=[:90]O)-[:-30]O-[:30]Et)-=-)}
\+
\chemfig{EtO-P(=[2]O)(-[6]O^{-})-OEt}
\schemestop
\par\vspace{6pt}

{\small A \omterm{def:b2:wittig:hwe}{reação de Horner--Wadsworth--Emmons} do fosfonoacetato de trietila com o benzaldeído. O
hidreto de sódio retira o hidrogênio situado entre o \omterm{def:b2:wittig:hwe}{fosfonato} e o éster; o ânion se adiciona ao
aldeído, e o intermediário de quatro membros se desfaz como na \omterm{def:b2:wittig:wittig}{reação de Wittig}. O produto é o
(\textit{E})-cinamato de etila; o subproduto é o ânion fosfato de dietila, solúvel em água.}
\end{omfigure}

\begin{proposition}[Por que os químicos gostam da reação HWE]\label{prop:b2:wittig:hwe}
A reação HWE dá sobretudo o alceno (\textit{E}), e seu subproduto fosforado, um sal solúvel em água,
é eliminado por uma simples lavagem aquosa; o ânion \omterm{def:b2:wittig:hwe}{fosfonato} é também mais nucleofílico que o \omterm{def:b2:wittig:ylide}{ilídeo
estabilizado} correspondente e reage também com as cetonas.
\end{proposition}

\begin{proof}[Raciocínio]
O ânion \omterm{def:b2:wittig:hwe}{fosfonato} é um carbânion estabilizado, de modo que sua adição é reversível e a reação ocorre
sob controle termodinâmico: (\textit{E}), como com os \omterm{def:b2:wittig:ylide}{ilídeos estabilizados}. Ao contrário de um \omterm{def:b2:wittig:ylide}{ilídeo
de fosfônio}, ele leva uma carga negativa inteira, não compensada por um cátion vizinho: ele é mais
reativo. O subproduto, \ce{(EtO)2PO2-}, é um sal iônico, enquanto o óxido de trifenilfosfina é um
sólido orgânico neutro que precisa ser separado do produto por cromatografia ou cristalização.
\end{proof}

\section{Escolher um método}

\begin{method}[Desconexão de Wittig]\label{met:b2:wittig:wittig-disconnection}
\begin{enumerate}
  \item Corte a ligação \ce{C=C} do alvo; um carbono torna-se um carbono carbonílico, o outro o carbono
        do ilídeo (um \omterm{def:b2:wittig:ylide}{sal de fosfônio}, a partir de um haleto).
  \item Das duas maneiras de fazê-lo, prefira aquela cujo haleto é primário (substituição bimolecular
        pela \ce{Ph3P}) e cujo composto carbonílico é um aldeído.
  \item Escolha o reagente pela geometria desejada: um ilídeo não estabilizado para (\textit{Z}), um
        \omterm{def:b2:wittig:ylide}{ilídeo estabilizado} ou um \omterm{def:b2:wittig:hwe}{fosfonato} (HWE) para (\textit{E}).
  \item Verifique que nada mais na molécula reage com a base usada.
\end{enumerate}
\end{method}

\begin{method}[Escolher uma reação que forma C=C]\label{met:b2:wittig:choose-cc}
\begin{center}
\small
\begin{tabular}{@{}lll@{}}
\hline
Método & Posição da \ce{C=C} & Geometria \\
\hline
Eliminação E2 ou E1 & regra de Zaitsev, misturas & sobretudo (\textit{E}), misturas \\
\omterm{def:b2:enolates-aldol:aldol}{Condensação aldólica} & conjugada com \ce{C=O} & (\textit{E}) \\
Wittig, ilídeo não estabilizado & na antiga \ce{C=O} & (\textit{Z}) \\
Wittig, \omterm{def:b2:wittig:ylide}{ilídeo estabilizado}; HWE & na antiga \ce{C=O} & (\textit{E}) \\
Alcino, depois catalisador envenenado & na antiga ligação tripla & (\textit{Z}) \\
\hline
\end{tabular}
\end{center}
\medskip
Para unir dois pedaços de tamanho semelhante em uma \ce{C=C}, a \omterm{def:b2:wittig:wittig}{reação de Wittig} e suas parentes são
em geral a primeira escolha; a metátese de olefinas, que troca as extremidades de dois alcenos, vem no
volume do 3.º ano.
\end{method}

\begin{history}[O ilídeo que se tornou reagente]
Georg Wittig, estudando os ilídeos de fósforo no início dos anos 1950, descobriu que o
metilidenotrifenilfosforano converte a benzofenona em 1,1-difenileteno e óxido de trifenilfosfina; com
seu aluno Ulrich Schöllkopf, ele transformou a observação em método geral em 1954. A indústria a adotou
depressa, entre outras para a vitamina A, e Wittig dividiu o Prêmio Nobel de Química de 1979.
\end{history}

\begin{safety}{\ghs[5mm]{GHS02}\,\ghs[5mm]{GHS05}\,\ghs[5mm]{GHS07}\,\ghs[5mm]{GHS08}}
As soluções de butil-lítio são inflamáveis e corrosivas e reagem violentamente com a água; o hidreto de
sódio libera hidrogênio com a água. A trifenilfosfina é nociva e corrosiva para os olhos, o fosfito de
trietila inflamável e nocivo; o óxido de trifenilfosfina é nocivo.
% ledger: ghs:BuLi, ghs:PPh3, ghs:triethyl-phosphite, ghs:Ph3PO
\end{safety}

\section{Exercícios}

\begin{exercise}[$\star$]\label{exo:b2:wittig:1}
Escreva a preparação do ilídeo \ce{Ph3P=CHCH3} a partir da trifenilfosfina e de um haloalcano.
\end{exercise}

\begin{exercise}[$\star$]\label{exo:b2:wittig:2}
Dê os produtos da \omterm{def:b2:wittig:wittig}{reação de Wittig} do benzaldeído com o ilídeo preparado a partir do iodometano.
\end{exercise}

\begin{exercise}[$\star$]\label{exo:b2:wittig:3}
Classifique como estabilizado, não estabilizado ou estabilizado por arila: \ce{Ph3P=CHCOOEt}, \ce{Ph3P=CHCH3},
\ce{Ph3P=CHCOCH3}, \ce{Ph3P=CHC6H5}. Qual deles exige butil-lítio?
\end{exercise}

\begin{exercise}[$\star$]\label{exo:b2:wittig:4}
Dê o produto do fosfonoacetato de trietila, do hidreto de sódio e do propanal, com sua geometria.
\end{exercise}

\begin{exercise}[$\star\star$]\label{exo:b2:wittig:5}
Deseja-se o metilidenociclo-hexano. Compare uma via de Wittig a partir da ciclo-hexanona com a
desidratação, catalisada por ácido, do 1-metilciclo-hexan-1-ol.
\end{exercise}

\begin{exercise}[$\star\star$]\label{exo:b2:wittig:6}
Desconecte o (\textit{E})-1,2-difenileteno (estilbeno) pela \omterm{def:b2:wittig:wittig}{reação de Wittig}. Por que a via de Wittig
é uma má escolha para o (\textit{E})-estilbeno puro, e o que você usaria em seu lugar?
\end{exercise}

\begin{exercise}[$\star\star$]\label{exo:b2:wittig:7}
O propanal reage com \ce{Ph3P=CHCH3} de um lado e com \ce{Ph3P=CHCOOEt} do outro. Dê os produtos e
sua geometria principal.
\end{exercise}

\begin{exercise}[$\star\star$]\label{exo:b2:wittig:8}
Calcule a economia atômica da metilenação da ciclo-hexanona por \ce{Ph3P=CH2}.
\end{exercise}

\begin{exercise}[$\star\star$]\label{exo:b2:wittig:9}
Escreva a reação de Michaelis--Arbuzov do fosfito de trietila com o bromoetanoato de etila, em duas
etapas: substituição pelo fósforo, depois desalquilação pelo brometo.
\end{exercise}

\begin{exercise}[$\star\star\star$]\label{exo:b2:wittig:10}
Proponha uma síntese do 1,6-difenil-hexa-1,3,5-trieno a partir do (\textit{E})-but-2-enodial
\ce{OHC-CH=CH-CHO} por duas \omterm{def:b2:wittig:wittig}{reações de Wittig}, e diga que ilídeo e em que quantidade.
\end{exercise}

\begin{exercise}[$\star\star\star$]\label{exo:b2:wittig:11}
A (\textit{E})-4-fenilbut-3-en-2-ona pode ser preparada por uma \omterm{def:b2:enolates-aldol:aldol}{condensação aldólica} ou por uma \omterm{def:b2:wittig:wittig}{reação
de Wittig}. Escreva as duas vias e compare-as.
\end{exercise}

\begin{exercise}[$\star\star\star$]\label{exo:b2:wittig:12}
Proponha duas vias para o (\textit{Z})-oct-4-eno, uma pela \omterm{def:b2:wittig:wittig}{reação de Wittig} e outra por meio de um
alcino, e diga que subprodutos cada uma deixa.
\end{exercise}

\section{Problema: Um feromônio por Wittig}

\begin{problem}[{Problema de fim de semana --- a retrossíntese de um alceno (Z), o ilídeo não
estabilizado e sua seletividade, o custo em óxido de trifenilfosfina, e a alternativa do alcino}]
\label{pb:b2:wittig:1}
Alvo: o (\textit{Z})-tricos-9-eno, \ce{CH3(CH2)7CH=CH(CH2)12CH3}, o feromônio da mosca doméstica da
introdução. Massas molares (\unit{g/mol}): C 12.011, H 1.008, O 15.999, P 30.974, Br 79.904.
% ledger: use:muscalure, aw:C, aw:H, aw:O, aw:P, aw:Br, ghs:nonanal, ghs:BuLi

\medskip
\noindent\textbf{Parte I --- Retrossíntese.}
\begin{enumerate}
  \item Escreva a fórmula molecular do alvo e verifique que é um alceno.
  \item Que ligação uma desconexão de Wittig corta?
  \item Dê os dois pares possíveis (composto carbonílico, \omterm{def:b2:wittig:ylide}{sal de fosfônio}).
  \item Para o par nonanal + ilídeo, que haloalcano fornece o \omterm{def:b2:wittig:ylide}{sal de fosfônio}?
  \item Por que um haloalcano primário convém para esta etapa?
  \item Escreva a formação do \omterm{def:b2:wittig:ylide}{sal de fosfônio}.
\end{enumerate}

\medskip
\noindent\textbf{Parte II --- O ilídeo e a geometria.}
\begin{enumerate}[resume]
  \item Que base desprotona esse sal? Por que o hidróxido de sódio não serviria?
  \item O ilídeo é estabilizado? Que geometria isso prevê?
  \item Escreva o \omterm{def:b2:wittig:wittig}{oxafosfetano} formado com o nonanal.
  \item Por que a reação deve ser conduzida em meio anidro e sob gás inerte?
  \item Por que a ligação dupla termina exatamente entre os carbonos 9 e 10?
  \item O que daria em vez disso a desidratação, catalisada por ácido, do tricosan-9-ol?
\end{enumerate}

\medskip
\noindent\textbf{Parte III --- O óxido de trifenilfosfina.}
\begin{enumerate}[resume]
  \item Escreva a equação global a partir do ilídeo e do nonanal.
  \item Calcule as massas molares do alvo, do nonanal, do \omterm{def:b2:wittig:ylide}{sal de fosfônio} e do
        óxido de trifenilfosfina.
  \item Calcule a economia atômica da etapa de Wittig (ilídeo + aldeído).
  \item Como se separa o óxido de trifenilfosfina de um hidrocarboneto?
  \item Por que a formação do óxido de trifenilfosfina é a força motriz?
  \item Por que a reação HWE não conviria a este alvo?
\end{enumerate}

\medskip
\noindent\textbf{Parte IV --- A via do alcino.}
\begin{enumerate}[resume]
  \item Que alcino, hidrogenado, dá o alvo, e com que catalisador?
  \item Construa esse alcino a partir do dec-1-ino e de um haloalcano: que haleto, que base?
  \item Por que a hidrogenação deve parar no alceno, e como o catalisador o garante?
  \item Compare as duas vias: número de etapas, controle da geometria, subprodutos.
  \item Qual é a economia atômica da etapa de hidrogenação?
  \item Dê a massa de óxido de trifenilfosfina coproduzida por grama de feromônio pela via de
        Wittig.
\end{enumerate}
\end{problem}
